\section{Introduction} \label{s:intro}

The subindex ev/odd is added to any space of functions on $\R$ to indicate its subspace of even/odd functions; in particular, $\Cinf=\Cinf_{\text{\rm ev}}\oplus\Cinf_{\text{\rm odd}}$ for $\Cinf:=\Cinf(\R)$. The Dunkl operator $T_\sigma$ ($\sigma>-1/2$) on $\Cinf$ is the perturbation of $\frac{{\rm d}}{{\rm d}x}$ def\/ined by $T_\sigma=\frac{{\rm d}}{{\rm d}x}$ on $\Cinf_{\text{\rm ev}}$ and $T_\sigma=\frac{{\rm d}}{{\rm d}x}+2\sigma\frac{1}{x}$ on $\Cinf_{\text{\rm odd}}$. The corresponding Dunkl harmonic oscillator is the perturbation $L_\sigma=-T_\sigma^2+s^2x^2$ of the harmonic oscillator $H=-\frac{{\rm d}^2}{{\rm d}x^2}+s^2x^2$ ($s>0$). The conjugation $E_\sigma=|x|^\sigma T_\sigma|x|^{-\sigma}$ on $|x|^\sigma\Cinf$ is equal to $\frac{{\rm d}}{{\rm d}x}-\sigma x^{-1}$ on $|x|^\sigma\Cinf_{\text{\rm ev}}$ and $\frac{{\rm d}}{{\rm d}x}+\sigma x^{-1}$ on $|x|^\sigma\Cinf_{\text{\rm odd}}$; note that $|x|^\sigma\Cinf_{\text{\rm ev/odd}}$ consists of even/odd functions, possibly not smooth or not even def\/ined at~$0$. Up to the product by a~constant, $E_\sigma$ was introduced by Yang~\cite{Yang1951}. In the form~$T_\sigma$, this operator was generalized to~$\R^n$ by Dunkl~\cite{Dunkl1988,Dunkl1989,Dunkl1991}, giving rise to what is now called Dunkl theory (see the survey~\cite{Rosler2003}); in particular, the Dunkl harmonic oscillator on~$\R^n$ was studied in \cite{Dunkl2002,NowakStempak2009b,NowakStempak2009a, Rosenblum1994}. See \cite{Plyushchay2000} for further generalizations on~$\R$. Sometimes the terms Yang--Dunkl operator and Yang--Dunkl harmonic oscillator are used in the case of~$\R$~\cite{Plyushchay2000}.

Let $p_k$ be the sequence of orthogonal polynomials for the measure $e^{-sx^2}|x|^{2\sigma}{\rm d}x$, taken with norm one and positive leading coef\/f\/icient. Up to normalization, these are the generalized Hermite polynomials~\cite[p.~380, Problem~25]{Szego1975}; see also \cite{Chihara1955,Chihara1978,DickinsonWarsi1963,DuttaChatterjeaMore1975,Rosenblum1994,Rosler1998}. The corresponding generalized Hermite functions are $\phi_k=p_ke^{-sx^2/2}$.

For each $m\in\N$, let $\SS^m$ be the Banach space of functions $\phi\in C^m(\R)$ with \mbox{$\sup_x|x^i\phi^{(j)}(x)|\!<\!\infty$} for $i+j\le m$; the corresponding Fr\'echet space $\SS=\bigcap_m\SS^m$ is the Schwartz space on $\R$. With domain $\SS$, $L_\sigma$ is essentially self-adjoint in $L^2(\R,|x|^{2\sigma}{\rm d}x)$, and the spectrum of its self-adjoint extension, $\LL_\sigma$, consists of the eigenvalues $(2k+1+2\sigma)s$ ($k\in\N$), with corresponding eigenfunctions $\phi_k$  \cite{Rosenblum1994}. For each real $m\ge0$, let $W_\sigma^m$ be the Hilbert space completion of $\SS$ with respect to the scalar product $\langle\phi,\psi\rangle_{W_\sigma^m}:=\langle(1+\LL_\sigma)^m\phi,\psi\rangle_\sigma$, where $\langle\ ,\ \rangle_\sigma$ denotes the scalar product of $L^2(\R,|x|^{2\sigma}{\rm d}x)$, obtaining a Fr\'echet space $W_\sigma^\infty=\bigcap_mW_\sigma^m$. We show the following embedding theorems; the second one is of Sobolev type.
